\section{Context: Mechanized Coinduction}
\label{sec:mechanized}

But\ldots why care about this when we have \code{cofix} in Rocq?

\subsection{Problems with Rocq's Native Coinduction}

\begin{enumerate}
  \item Non-compositional
  \item Easy to make mistakes with \code{cofix}
  \item Guess the bisimulation candidate \emph{up front}
  \item Not amenable to automation
\end{enumerate}

Another (big) problem:
\begin{enumerate}\setcounter{enumi}{4}
  \item No support for \textbf{up-to} techniques
\end{enumerate}

\subsection{Coinduction Up-to}

Motivation: easier proof by coinduction.

\code{R <= b R} $\longrightarrow$ \code{R <= b (f R)} for some \code{f}

Example: \code{f x := x \/ id}. This is \textbf{``up to reflexivity''}. Could
prove manually, but\ldots would rather show \code{~} is reflexive!

Other common up-to functions are symmetry, transitivity, and up-to context.
Some up-to techniques, like up-to context, are \underline{essential} for proofs
in sophisticated developments.

\subsection{(Mechanized) Coinduction Wishlist}
\label{sec:wishlist}

\begin{enumerate}
  \item Compositional (Non-syntactic guardedness)
  \item Smooth way to prove soundness of \textbf{up-to techniques}
  \item Straightforward to use
  \item Amenable to automation
  \item Nice-to-have: incremental bisimulation
\end{enumerate}

How about just formalizing Knaster-Tarski? \gap{The scorecard from talk slide 35:
1~\checkmark, 2~$\times$, 3~\checkmark, 4~\checkmark, 5~$\times$.}

\note{\textbf{This wishlist is the paper's spine and should be presented as
such.} The talk scores it four times, at slides 35 (Knaster-Tarski), 41 (paco),
64 (tower) and 89 (the rewrite, items 6--10). That repetition is what makes the
talk's argument land, and it is a structure a paper can keep cheaply as one
table with four columns rather than four separate scorecards. Worth deciding
early, because it also fixes how \Cref{sec:results} ends.}

\subsection{The Pursuit of Happiness in Coinduction}

\gap{The timeline from talk slide 37: \code{cofix} $\to$ \code{paco} $\to$ the
companion (Pous~\cite{pous-companion}; Parrow \& Weber; Schäfer \&
Smolka~\cite{schafer-smolka}) $\to$ ITrees and
\code{gpaco} $\to$ New ITrees. This is the one context figure that earns a full
column in the paper, because it places the two libraries the paper compares.}
